% Drawing 141/01-5 -- folder 141, batch 01, archive page 10 (PDF page 11).
% Transcription: au-dessous, une ligne horizontale portant sept points marqués, prolongée à gauche par deux points et à droite par trois points de suspension (le cas $n = \infty$).
% Kind: pointed line continued by dots on both sides (n infinite). Method 2 (schematic redraw from the facsimile). Status: ok.
% Notes: As on the page: a horizontal line with seven marked points (its two ends and five between, unevenly spaced, each computed on the line), two dots before it at the left and three after it at the right. A stray dot under the line, near the middle, is left out. Redrawn 2026-11-12 (issue 128): to check again.
% Coordinates: 1 unit = 1 mm = 1/100 of the crop width (bbox in page fractions: [0.13, 0.385, 0.56, 0.425]).
% Redrawn by Claude (claude-fable-5-2), issue #128, 2026-11-12; not checked by a human; not compiled (no TeX in the session).
% Rights: reproduces part of the work -- behind RENDER_TRANSCRIPTS, never published by itself.
\documentclass[tikz,border=3pt]{standalone}
\usepackage{amsmath,amssymb}
\usetikzlibrary{patterns,arrows.meta}
\tikzset{pencil/.style={black!45}, struck/.style={black!35}, mpt/.style={circle,draw,inner sep=1.3pt}, pt/.style={circle,fill,inner sep=0.9pt}, opt/.style={circle,draw,fill=white,inner sep=1.1pt}}
\begin{document}
\begin{tikzpicture}[x=1mm,y=1mm,line cap=round,line join=round,font=\scriptsize]
\draw (13.5,-8.8) -- (83,-10.5);
\node[pt] at (13.5,-8.8) {};
\node[pt] at (25,-9.1) {};
\node[pt] at (40,-9.4) {};
\node[pt] at (52,-9.7) {};
\node[pt] at (61.5,-10) {};
\node[pt] at (73.5,-10.3) {};
\node[pt] at (83,-10.5) {};
\node[pt] at (5.5,-8.5) {};
\node[pt] at (9,-8.6) {};
\node[pt] at (87,-10.8) {};
\node[pt] at (89.3,-11) {};
\node[pt] at (91.5,-11.2) {};
\end{tikzpicture}
\end{document}
